\documentclass[11pt,a4paper]{article}
\usepackage[margin=2.1cm]{geometry}
\usepackage{booktabs}
\usepackage{amsmath}
\usepackage{amssymb}
\usepackage{enumitem}
\usepackage{xcolor}
\usepackage[breakable]{tcolorbox}
\usepackage{fancyhdr}
\usepackage{titlesec}
\usepackage[colorlinks=true,linkcolor=blue!45!black,urlcolor=blue!45!black]{hyperref}

\definecolor{navy}{RGB}{31,56,100}
\definecolor{okgreen}{RGB}{0,97,0}
\definecolor{warnred}{RGB}{156,0,6}
\definecolor{amber}{RGB}{156,101,0}
\definecolor{barblue}{RGB}{60,110,180}

\newcommand{\good}[1]{\textcolor{okgreen}{\textbf{#1}}}
\newcommand{\bad}[1]{\textcolor{warnred}{\textbf{#1}}}
\newcommand{\neut}[1]{\textcolor{amber}{\textbf{#1}}}
\newcommand{\histbar}[1]{\textcolor{barblue}{\rule{#1}{7pt}}}

\titleformat{\section}{\Large\bfseries\color{navy}}{\thesection}{0.6em}{}
\titleformat{\subsection}{\large\bfseries\color{navy}}{\thesubsection}{0.6em}{}

\newtcolorbox{say}[1]{colback=blue!3,colframe=navy,fonttitle=\bfseries,
  title={SAY --- #1},breakable,left=4pt,right=4pt,top=4pt,bottom=4pt}
\newtcolorbox{ifasked}[1]{colback=gray!6,colframe=gray!55,fonttitle=\bfseries,
  title={IF ASKED --- #1},breakable,left=4pt,right=4pt,top=4pt,bottom=4pt}

\pagestyle{fancy}
\fancyhf{}
\fancyhead[L]{\small\color{navy}Log Report \#006 --- Benchmark Artifact Evaluation}
\fancyhead[R]{\small\thepage}
\renewcommand{\headrulewidth}{0.4pt}

\title{\vspace{-1.4cm}\textbf{Log Report \#006}\\[2mm]
\large Benchmark Artifact Evaluation: Measuring the Ceiling on Real Instances\\[1mm]
\normalsize GPU-Accelerated Nemhauser--Trotter Kernelization for WCSPs}
\author{\normalsize Presentation notes --- read alongside \texttt{artifact\_results.xlsx}}
\date{\normalsize 2 September 2026}

\begin{document}
\maketitle
\thispagestyle{fancy}

\begin{tcolorbox}[colback=okgreen!8,colframe=okgreen,title=\textbf{THE RESULT}]
For every instance in the published benchmark we measured the \emph{maximum} that any minimum-cut
kernelizer could possibly resolve --- a bound that constrains Gurobi exactly as much as it
constrains us. On the UAI collection our kernelizer already attains that maximum on
\textbf{145 of 146} instances (99.3\,\%), and on \textbf{97} of them the maximum is \textbf{zero}.
On the \texttt{evalgm} collection we attain it on \textbf{348 of 373} (93.3\,\%).

The 25 instances that fall short on \texttt{evalgm} are \emph{exactly} the 24 where max-flow trails
Gurobi, plus one more --- the diagnostic predicts the gap with no error. Those 25 carry integral
weights, so weight perturbation is permitted on them, unlike on UAI. Three of them are structurally
identical to the synthetic \texttt{d\_m2000} case that perturbation already takes from 0\,\% to
100\,\%.
\end{tcolorbox}

\vspace{2mm}
\noindent\textbf{How to use this document.} Keep it open on a second screen; show the tables from
the workbook and read from the shaded boxes.

\tableofcontents

% ======================================================================
\section{What this run set out to answer}

\begin{say}{opening}
Everything until now was measured on instances we generated ourselves. The obvious objection is
that we chose problems that flatter our method. So we obtained the published benchmark artifact
from the paper we build on, and evaluated on it.

We did not simply run the kernelizer and report numbers. We first measured, for every single
instance, the \emph{maximum any method of this kind could achieve}. That bound applies to Gurobi
too. Only then does a comparison mean anything.
\end{say}

The instrument is \texttt{apps/mincut\_lattice.cpp}. For each variable it asks whether the two
copies in the double cover lie in the same strongly connected component of the residual graph. Two
nodes in one component are reachable from each other along edges that still have spare capacity, so
a cut separating them would leave a residual edge crossing it --- which no minimum cut can do. Hence:

\begin{itemize}[leftmargin=1.4em,itemsep=1pt]
  \item \textbf{same component} $\Rightarrow$ same side in \emph{every} minimum cut $\Rightarrow$
        undecidable by \emph{any} method of this kind, Gurobi included;
  \item \textbf{different components} $\Rightarrow$ some minimum cut separates them.
\end{itemize}

Counting the second group gives the \textbf{ceiling}.

% ======================================================================
\section{What was attempted --- show Table 1}

\begin{center}
\begin{tabular}{@{}lrrrrl@{}}
\toprule
Collection & Files & Header-Boolean & Loadable & Completed & Note \\
\midrule
\texttt{uai-instances} & 295 & 160 & 160 & 146 & 14 timed out at 900\,s \\
\texttt{evalgm}, $<$10k vars & 557 & 557 & \bad{373} & 373 & 184 rejected by the loader \\
\bottomrule
\end{tabular}
\end{center}

\begin{say}{Table 1, and a correction we found}
Two collections. The UAI competition instances, and the smaller half of the \texttt{evalgm}
benchmark.

One correction. Our earlier census counted usable instances from the file header, which declares a
maximum domain size. But the loader also checks \emph{each} variable, and rejects any with domain
size one --- a variable pinned to a single value. On this subset that removed a further 184 of 557.
So the figure of 878 usable \texttt{evalgm} instances is an over-count and we will re-measure it
properly.
\end{say}

% ======================================================================
\section{UAI: the ceiling is almost always zero --- show Table 2}

\begin{center}
\begin{tabular}{@{}lrrl@{}}
\toprule
Verdict & Instances & Share & Meaning \\
\midrule
\bad{STRUCTURAL} & \textbf{97} & 60.6\,\% & ceiling is exactly zero \\
\good{OPTIMAL}   & \textbf{48} & 30.0\,\% & we already attain the ceiling \\
\neut{SELECTION} & 1 & 0.6\,\% & a better cut exists (\texttt{relational\_1}) \\
\neut{TIMEOUT}   & 14 & 8.8\,\% & unfinished at 900\,s, all CNF-derived \\
\bottomrule
\end{tabular}
\end{center}

\begin{center}
\begin{tabular}{@{}rl r@{}}
\toprule
Ceiling & & Instances \\
\midrule
exactly 0\,\% & \histbar{8.2cm} & \textbf{97} \\
0 -- 1\,\%    & \histbar{2.3cm} & 27 \\
1 -- 10\,\%   & \histbar{0.34cm} & 4 \\
10 -- 50\,\%  & & 0 \\
50 -- 99\,\%  & & 0 \\
100\,\%       & \histbar{1.5cm} & \textbf{18} \\
\bottomrule
\end{tabular}
\end{center}

\begin{say}{Table 2}
On 145 of the 146 instances that finished, our kernelizer already achieves the theoretical maximum
--- ninety-nine point three percent. Not close to Gurobi; \emph{at the ceiling}, which Gurobi cannot
exceed either.

And on ninety-seven of them the ceiling is exactly zero. Nothing can be resolved by any method of
this kind. That is not our failure; it is a property of those problems.

Look at the distribution. Nothing lands between ten percent and one hundred. The \emph{ceiling
itself} is all-or-nothing, which means the all-or-nothing behaviour is a property of the instances
rather than of any algorithm. That fully explains the histogram published in the original paper.
\end{say}

\begin{ifasked}{what about the one SELECTION instance?}
\texttt{relational\_1}, ceiling 100\,\%, currently 0\,\%. It is also the one instance our
perturbation improved in the end-to-end run, 50\,\% to 100\,\%. So perturbation did little on UAI
not because it fails, but because there was exactly one thing it could fix, and it fixed it.
\end{ifasked}

\begin{ifasked}{why did perturbation not run on the rest?}
It refuses to. UAI weights come from log-probabilities and are decimals. The safety argument needs
integral weights: two solutions of different cost then differ by at least one, which the
perturbation can never bridge. With decimals that is false and the answer could change. The
guard is now explicit and reports \texttt{DISABLED: vertex weights are not integral}.
\end{ifasked}

% ======================================================================
\section{evalgm: a real target appears --- show Table 3}

\begin{center}
\begin{tabular}{@{}lrrl@{}}
\toprule
Verdict & Instances & Share of loaded & Meaning \\
\midrule
\bad{STRUCTURAL} & 222 & 59.5\,\% & ceiling is exactly zero \\
\good{OPTIMAL}   & 126 & 33.8\,\% & already at the ceiling \\
\neut{SELECTION} & \textbf{25} & \textbf{6.7\,\%} & \good{fixable, and perturbation is permitted} \\
SKIP             & 184 & --- & rejected by the loader \\
\bottomrule
\end{tabular}
\end{center}

\begin{center}
\begin{tabular}{@{}rl r@{}}
\toprule
Ceiling & & Instances \\
\midrule
exactly 0\,\% & \histbar{7.4cm} & \textbf{222} \\
0 -- 1\,\%    & \histbar{0.9cm} & 27 \\
1 -- 10\,\%   & \histbar{0.27cm} & 8 \\
10 -- 50\,\%  & & 0 \\
50 -- 99\,\%  & \histbar{2.1cm} & \textbf{63} \\
100\,\%       & \histbar{1.8cm} & \textbf{53} \\
\bottomrule
\end{tabular}
\end{center}

\begin{say}{Table 3}
The same shape, but with an important difference: a hundred and sixteen instances have a ceiling
above fifty percent. Real headroom exists here, and we already reach it on three hundred and
forty-eight of three hundred and seventy-three --- ninety-three point three percent.

Crucially, every one of the three hundred and seventy-three completed instances has \emph{integral}
weights. That is the condition perturbation needs, and it holds throughout this collection.
\end{say}

% ======================================================================
\section{The diagnostic predicts the gap exactly --- show Table 6}

\begin{center}
\begin{tabular}{@{}lrrrrr@{}}
\toprule
Collection & Comparable & Identical & Max-flow better & Max-flow worse & Agreement \\
\midrule
UAI (Boolean) & 141 & 135 & 5 & 1 & \good{95.7\,\%} \\
\texttt{evalgm} ($<$10k) & 373 & 349 & 0 & \textbf{24} & \good{93.6\,\%} \\
\bottomrule
\end{tabular}
\end{center}

\begin{say}{the sharpest finding in this report}
On \texttt{evalgm} our kernelizer matches Gurobi exactly on three hundred and forty-nine of three
hundred and seventy-three instances, and never beats it. It loses on twenty-four.

The ceiling analysis, run completely independently, flagged twenty-five instances as having room to
improve. Those twenty-five \emph{contain} the twenty-four. The diagnostic identifies precisely where
we lose, without ever consulting Gurobi.

That is what makes the tool worth reporting in its own right: it tells you, cheaply and per
instance, whether an instance is worth working on.
\end{say}

\subsection{Against the published results}

\begin{center}
\begin{tabular}{@{}lr@{}}
\toprule
Comparable instances & 150 \\
Exact match & \good{125} \\
Differ & 25 \\
Agreement & \good{83.3\,\%} \\
\bottomrule
\end{tabular}
\end{center}

\begin{say}{on the differences}
Variable counts matched the published figure on \emph{all} one hundred and fifty-five instances,
including every one where the resolved count differs. So our parser reads the same problem the
authors did. The differences are in both directions and are concentrated on planning instances,
which is the signature of LP tie-breaking between different Gurobi versions --- the published runs
used a 2017 build.
\end{say}

% ======================================================================
\section{The 25 fixable instances --- show Table 5}

\begin{center}
\small
\begin{tabular}{@{}llrrrr@{}}
\toprule
Instance & Family & Vertices & Now & Ceiling & Recoverable \\
\midrule
\texttt{hamming10-2.clq} & MaxClique & 11{,}264 & \bad{0\,\%} & \good{100\,\%} & \textbf{11{,}264} \\
\texttt{hamming8-2.clq}  & MaxClique &  2{,}304 & \bad{0\,\%} & \good{100\,\%} & \textbf{2{,}304} \\
\texttt{hamming6-2.clq}  & MaxClique &    448 & \bad{0\,\%} & \good{100\,\%} & \textbf{448} \\
\texttt{cat\_paths\_60\_150\_0001} & Auction & 4{,}770 & 0\,\% & 2.12\,\% & 101 \\
\texttt{cat\_paths\_60\_170\_0007} & Auction & 5{,}469 & 0\,\% & 0.84\,\% & 46 \\
\texttt{cat\_paths\_60\_140\_0006} & Auction & 4{,}205 & 0\,\% & 0.76\,\% & 32 \\
\multicolumn{4}{l}{\emph{\ldots and 17 further Auction / PackupWeighted instances}} & & 267 \\
\midrule
\multicolumn{5}{r}{\textbf{Total recoverable}} & \textbf{14{,}462} \\
\bottomrule
\end{tabular}
\end{center}

\begin{say}{Table 5 --- the next experiment}
Twenty-five instances, fourteen thousand four hundred and sixty-two recoverable vertices. Ninety-seven
percent of that sits in three instances from the maximum-clique family.

Those three are worth looking at closely. All-half optimal, integral weights, ceiling one hundred
percent, currently resolving zero. That is exactly the profile of \texttt{d\_m2000} --- the synthetic
instance where perturbation already takes us from zero to one hundred percent. If perturbation
behaves the same way here, it will be on a third-party benchmark instance rather than one we made.
\end{say}

\begin{center}
\begin{tabular}{@{}lrrlll@{}}
\toprule
Instance & Vertices & $W$ & Integral & All-half optimal & Ceiling \\
\midrule
\texttt{hamming6-2.clq}  &    448 & 25{,}024 & yes & YES & 100\,\% \\
\texttt{hamming8-2.clq}  &  2{,}304 & 526{,}592 & yes & YES & 100\,\% \\
\texttt{hamming10-2.clq} & 11{,}264 & $1.05\times10^{7}$ & yes & YES & 100\,\% \\
\midrule
\texttt{d\_m2000} (synthetic, for comparison) & 4{,}430 & 7{,}016 & yes & YES & 100\,\% \\
\bottomrule
\end{tabular}
\end{center}

% ======================================================================
\section{Family matters more than size --- show Table 4}

\begin{center}
\begin{tabular}{@{}lrrrrrr@{}}
\toprule
Family & Structural & Optimal & Selection & Skipped & Total & At ceiling \\
\midrule
\texttt{CFN/Auction}          & 144 &  6 & \neut{20} &   0 & 170 & 88.2\,\% \\
\texttt{MRF/DBN}              &   0 &  0 &  0 & 108 & 108 & --- \\
\texttt{CVPR/ChineseChars}    &   0 & 63 &  0 &   0 &  63 & \good{100\,\%} \\
\texttt{WPMS/MaxClique}       &  59 &  0 & \neut{3} &   0 &  62 & 95.2\,\% \\
\texttt{WPMS/PackupWeighted}  &   1 &  2 & \neut{2} &  57 &  62 & 60.0\,\% \\
\texttt{MRF/Segmentation}     &   0 & 50 &  0 &   0 &  50 & \good{100\,\%} \\
\texttt{MRF/Grid}             &  16 &  5 &  0 &   0 &  21 & \good{100\,\%} \\
\texttt{WPMS/PlanningWithPref}&   0 &  0 &  0 &  13 &  13 & --- \\
\texttt{WPMS/MIPLib}          &   0 &  0 &  0 &   6 &   6 & --- \\
\texttt{MaxCSP/Coloring}      &   2 &  0 &  0 &   0 &   2 & \good{100\,\%} \\
\bottomrule
\end{tabular}
\end{center}

\begin{say}{Table 4}
Behaviour clusters by family, not by size. Segmentation and Chinese-character instances are at the
ceiling everywhere. Auction and MaxClique hold every fixable instance we found. Four families are
lost entirely to the loader's domain restriction.
\end{say}

% ======================================================================
\section{Corrections to Report 005 --- show Table 7}

\begin{enumerate}[leftmargin=1.6em]
\item \textbf{Usable \texttt{evalgm} count.} 878 was derived from file headers. The loader also
      rejects variables of domain size 1, removing a further 184 of 557 on this subset. The true
      figure is lower and must be re-measured with a loader-accurate test.
\item \textbf{The all-half law is one-way.} ``max-flow $=W$ $\Rightarrow$ resolves zero'' holds
      everywhere. The converse fails on UAI: 35 instances have all-half not optimal yet a ceiling of
      zero. On \texttt{evalgm} both directions hold on all 373. The two-way form appears to require
      integral weights.
\item \textbf{Perturbation is not applicable to UAI} by design, not by failure.
\item \textbf{The sweep's verdict counts were mislabelled.} 149 reported as \textsc{timeout} were in
      fact 135 non-Boolean skips and 14 real timeouts. The script label is fixed.
\end{enumerate}

% ======================================================================
\section{What is good, and what is not}

\subsection*{Good}
\begin{itemize}[leftmargin=1.4em]
\item 99.3\,\% at the ceiling on UAI, 93.3\,\% on \texttt{evalgm}, on third-party instances.
\item The ceiling bounds Gurobi too, so ``we cannot do better'' becomes ``nobody can''.
\item The diagnostic predicts the Gurobi gap exactly: 25 flagged, 24 real losses, all contained.
\item 83.3\,\% exact agreement with the published per-instance results; variable counts match on
      all 155, so the parser is correct.
\item A clear next target: 25 instances, integral weights, perturbation permitted, and three of
      them structurally identical to a case perturbation already solves.
\end{itemize}

\subsection*{Not good}
\begin{itemize}[leftmargin=1.4em]
\item 14 UAI timeouts at 900\,s, all CNF-derived, where Gurobi finishes in seconds. Max-flow is
      genuinely slower on that instance class.
\item Max-flow never beats Gurobi on \texttt{evalgm} --- 349 ties, 24 losses, 0 wins.
\item The usable-instance count was over-stated and must be corrected.
\item \texttt{logistics.c.cnf} recorded 5{,}240\,s against a 900\,s limit, most likely a machine
      sleep. That instance should be re-run before the number is quoted.
\item Only the smaller half of \texttt{evalgm} has been evaluated; instances above 10{,}000
      variables remain.
\end{itemize}

% ======================================================================
\section{Next steps}

\begin{enumerate}[leftmargin=1.6em]
\item Run perturbation on the three \texttt{hamming} instances. Prediction: 0\,\% $\to$ 100\,\%,
      matching \texttt{d\_m2000}.
\item Run perturbation across all 25 SELECTION instances and verify the objective is unchanged with
      an exact solve.
\item Re-count usable \texttt{evalgm} instances with a loader-accurate test.
\item Extend the sweep to \texttt{evalgm} instances above 10{,}000 variables.
\item Investigate why max-flow times out on CNF-derived instances.
\item Re-run \texttt{logistics.c.cnf} under a controlled timer.
\end{enumerate}

\vspace{4mm}
\noindent\rule{\textwidth}{0.4pt}\\
\small
Data in \texttt{artifact\_results.xlsx}. Raw logs in \texttt{log/2026-08-31/},
\texttt{log/2026-09-01/}, \texttt{log/2026-09-02/}. Branch \texttt{study/artifact-eval}.
Hardware: Intel i5-14450HX, 22.6\,GiB RAM, RTX 4050 Laptop (6\,GiB, sm\_89), Ubuntu 26.04,
g++ 15.2.0, CUDA 12.4, Gurobi 13.0.2 single-thread.

\end{document}
